\section{Comparison with EigenLayer}

\label{sec:comparison}
%%==========================================================================

Table~\ref{tab:eigenlayer} provides a structured comparison between LRP and EigenLayer.

\begin{table}[H]
\centering
\caption{LRP vs.\ EigenLayer Structural Comparison}
\label{tab:eigenlayer}
\begin{tabular}{p{3.5cm}p{4.5cm}p{4.5cm}}
\toprule
\textbf{Dimension} & \textbf{LRP (Lux)} & \textbf{EigenLayer (Ethereum)} \\
\midrule
Native finality & Sub-second (Beacon) & 12-second slots + finality epochs \\
Slashing latency & $< 2$s end-to-end & 7-day challenge period minimum \\
Slashing enforcement & P-Chain native + C-Chain EVM & Ethereum EVM only \\
Operator registration & BLS attestation, P-Chain native & Beacon chain withdrawal creds \\
Capital form & Native LUX or delegated LUX & ETH, stETH, rETH, LSTs \\
Unbonding period & 7--90 days (AVS-configurable) & 7 days (Ethereum standard) \\
Multi-appchain support & Native (P-Chain appchain registry) & External (AVS-level coordination) \\
Gas cost (slash tx) & $\approx\$0.01$ & $\approx\$15$--$\$80$ \\
Governance & LUX token governance & EIGEN token governance \\
Insurance fund & Protocol-native, 30\% of slashes & None (AVS-level only) \\
\bottomrule
\end{tabular}
\end{table}

\paragraph{Structural Advantages of LRP.}
We identify five structural advantages arising from Lux's architecture:

\begin{enumerate}[leftmargin=1.5em]
  \item \textbf{Native appchain integration}: Lux appchains are first-class citizens of the P-Chain.
    LRP can leverage P-Chain state transitions for slashing without cross-chain messaging,
    eliminating a major latency and gas cost bottleneck.
  \item \textbf{Sub-second slashing}: Beacon finality allows slashing decisions to be
    finalized in under 2 seconds, versus Ethereum's minimum 7-day window for LST-based restaking.
  \item \textbf{Lower costs}: P-Chain transactions cost $< \$0.01$, versus \$15--\$80 for
    Ethereum slashing coordination.
  \item \textbf{Explicit allocation model}: LRP's $\alpha_{v,s}$ allocation fractions provide
    explicit security accounting; EigenLayer uses implicit allocation through shared collateral
    without per-AVS fractions.
  \item \textbf{Protocol-native insurance}: The LRP insurance fund is enforced at the protocol
    layer; EigenLayer relies on AVS-level indemnity agreements with no protocol backstop.
\end{enumerate}

% Added 2025-04: Bitcoin restaking integration
\subsection{Bitcoin Restaking via Babylon}
\label{sec:babylon}

Babylon~\cite{babylon2023bitcoin} provides Bitcoin-native staking by locking BTC in
self-custodial UTXO covenants with time-locked slashing scripts.  A critical distinction:
Babylon's BTC staking prevents \emph{long-range attacks} on PoS chains by anchoring
finality to Bitcoin's proof-of-work, but it does \textbf{not} provide Byzantine fault
tolerance in execution environments.  BFT requires validators to run the chain's software
and attest to state transitions; Babylon stakers do neither.  The two mechanisms are
therefore complementary, not substitutive.

Lux integrates Babylon through two components:

\paragraph{On-chain: \texttt{BabylonBTCStrategy}.}
A C-Chain strategy contract registered in the AVS marketplace that accepts LBTC (wrapped
BTC on Lux) and stakes it into Babylon's finality-provider set.  The contract exposes:
\begin{itemize}[leftmargin=1.5em]
  \item \texttt{deposit(uint256 amount)}: locks LBTC and delegates to a Babylon finality
    provider;
  \item \texttt{withdraw(uint256 amount)}: initiates Babylon unbonding ($\approx$10 days);
  \item \texttt{slash(bytes proof)}: forwards Babylon slashing evidence to the LRP
    coordinator.
\end{itemize}

\paragraph{Cross-chain: Teleporter bridge.}
BTC locked on Bitcoin L1 is attested by an MPC watcher network (threshold $t$-of-$n$,
$n \geq 15$, $t \geq 10$).  Upon confirmation, LBTC is minted on Lux via the Teleporter
protocol~\cite{lux2025teleporter}.  Yield accrued by the Babylon finality provider flows
back to Lux through a \texttt{mintYield()} callback on the Teleporter bridge contract,
credited to the \texttt{BabylonBTCStrategy} vault.

\paragraph{Security accounting.}
BTC restaked via Babylon contributes to the economic security of Lux appchains at the
BTC/LUX oracle price, subject to a 20\% haircut to account for bridge and oracle risk.
For an appchain $s$:
\[
  \Pi_{\text{total}}(s) = \Pi_{\text{LUX}}(s) + 0.8 \cdot \Pi_{\text{BTC}}(s).
\]

This extends the restaking hypergraph: BTC-backed commitments form a separate edge set
$\mathcal{E}_{\text{BTC}} \subset \mathcal{E}$ with their own slashing conditions
(Babylon covenant violations) distinct from LUX-native slashing.

% Added 2025-04: Cross-chain restaking economics
\subsection{Cross-Chain Restaking Economics}
\label{sec:crosschain}

Bridged assets beyond BTC---including ETH and SOL---can participate in Lux's restaking
ecosystem.  Each asset class maps to its native yield source:

\begin{table}[H]
\centering
\caption{Cross-Chain Restaking Asset Strategies}
\label{tab:crosschain}
\begin{tabular}{llll}
\toprule
\textbf{Asset} & \textbf{Bridge} & \textbf{Native Yield Source} & \textbf{Lux Strategy} \\
\midrule
ETH   & Teleporter (MPC) & EigenLayer restaking & \texttt{EigenETHStrategy} \\
BTC   & Teleporter (MPC) & Babylon finality provision & \texttt{BabylonBTCStrategy} \\
SOL   & Wormhole NTT     & Jito/Marinade liquid staking & \texttt{SolanaLSTStrategy} \\
LUX   & Native           & LRP validator staking & \texttt{NativeLUXStrategy} \\
\bottomrule
\end{tabular}
\end{table}

The xLUX receipt token aggregates yield across all strategies.  When an operator opts into
an AVS, their vault may hold a mix of asset-specific strategy positions.  xLUX
represents a pro-rata claim on the combined yield:
\[
  \text{xLUX}_{\text{yield}} = \sum_{a \in \{\text{LUX, ETH, BTC, SOL}\}}
    w_a \cdot r_a,
\]
where $w_a$ is the portfolio weight of asset $a$ and $r_a$ its epoch yield rate.  This
creates a unified yield layer: operators and delegators hold a single token whose return
reflects the aggregate economics of multi-chain restaking.

Multi-chain restaking increases total economic security backing Lux.  If $W_a$ denotes the
total value of asset $a$ restaked, aggregate security becomes:
\[
  \Pi_{\text{multi}}(s) = \sum_{a} \gamma_a \cdot W_a \cdot \bar{\alpha}_s,
\]
where $\gamma_a \in (0,1]$ is the haircut factor for asset $a$ (1.0 for native LUX,
0.8 for BTC, 0.85 for ETH, 0.75 for SOL) and $\bar{\alpha}_s$ is the mean allocation
to service $s$.

% Added 2025-04: Restaking protocol comparison
\subsection{Multi-Protocol Comparison}
\label{sec:multicomparison}

Table~\ref{tab:multicompare} extends the EigenLayer comparison in
Table~\ref{tab:eigenlayer} to include Babylon and Symbiotic~\cite{symbiotic2024protocol}.

\begin{table}[H]
\centering
\caption{Restaking Protocol Comparison: LRP vs.\ EigenLayer vs.\ Babylon vs.\ Symbiotic}
\label{tab:multicompare}
\begin{tabular}{p{2.8cm}p{2.6cm}p{2.6cm}p{2.6cm}p{2.6cm}}
\toprule
\textbf{Dimension} & \textbf{LRP (Lux)} & \textbf{EigenLayer} & \textbf{Babylon} & \textbf{Symbiotic} \\
\midrule
Base asset & LUX + bridged & ETH + LSTs & BTC (native) & Any ERC-20 \\
Security type & BFT + economic & BFT + economic & Long-range attack prev. & Economic only \\
Slashing & On-chain, $<$2s & On-chain, 7d & Bitcoin covenant & Vault-configurable \\
Finality & Sub-second & 12s slots & Bitcoin-anchored & Inherits L1 \\
Operator model & Validator opt-in & Delegation & Finality providers & Vault curators \\
Appchain native & Yes (P-Chain) & No & No & No \\
Cross-chain & Teleporter & EVM bridges & BTC $\to$ PoS & EVM only \\
Yield token & xLUX & None native & None native & None native \\
Insurance & Protocol fund & AVS-level & None & Vault-level \\
\bottomrule
\end{tabular}
\end{table}

Key distinctions:

\begin{enumerate}[leftmargin=1.5em]
  \item \textbf{Babylon is not BFT.}  Babylon stakers do not validate state transitions.
    They provide economic finality anchored to Bitcoin's PoW, preventing long-range attacks
    but not equivocation or invalid execution.  LRP provides both BFT (through validator
    restaking) and can layer Babylon's long-range protection on top via
    \texttt{BabylonBTCStrategy}.
  \item \textbf{Symbiotic is vault-generic, not chain-native.}  Symbiotic allows arbitrary
    ERC-20 collateral with configurable slashing, but lacks native appchain integration or
    sub-second finality.  Its permissionless vault model trades protocol-level safety
    guarantees for flexibility.
  \item \textbf{EigenLayer is ETH-bound.}  EigenLayer's deep integration with Ethereum's
    beacon chain gives it native ETH restaking but limits cross-chain participation.
    LRP accepts multi-asset collateral natively.
  \item \textbf{LRP unifies yield.}  xLUX is the only receipt token that aggregates yield
    across multiple base assets and their respective native yield strategies, creating a
    single instrument for diversified restaking exposure.
\end{enumerate}

%%==========================================================================
